\chapter{Persistencia políglota}
\section{PostgreSQL}
PostgreSQL respalda Identity y Catalog. Identity mantiene \texttt{users}; Catalog emplea \texttt{movie}, \texttt{genre}, \texttt{artist}, tablas de relación, fuentes de video y subtítulos. Las relaciones N:M preservan la consistencia entre películas, géneros y artistas, y Flyway versiona el esquema.
\begin{figure}[H]\centering\includegraphics[width=.92\textwidth]{figures/er-postgresql.pdf}\caption{Modelo relacional PostgreSQL para Identity y Catalog}\label{fig:er_postgresql}\end{figure}
\section{MySQL}
Community utiliza MySQL para los modelos de usuarios, clubes, membresías, salas y participantes. La tabla \texttt{memberships} materializa la relación usuario--club; las salas guardan anfitrión, película y club opcional. El identificador de película es una referencia de aplicación al catálogo, no una clave foránea interservicio.
\begin{figure}[H]\centering\includegraphics[width=.92\textwidth]{figures/er-mysql.pdf}\caption{Modelo relacional MySQL del Community Service}\label{fig:er_mysql}\end{figure}
\section{MongoDB}
Interaction usa MongoDB para \texttt{reviews}, \texttt{likes}, \texttt{watchlists} y \texttt{watch\_history}. Cada documento combina la acción con snapshots de usuario y película para facilitar la lectura contextual y la extracción analítica. Los índices compuestos únicos de reseñas y likes evitan duplicados por usuario y película.
\begin{figure}[H]\centering\includegraphics[width=.92\textwidth]{figures/modelo-mongodb.pdf}\caption{Modelo documental MongoDB del Interaction Service}\label{fig:modelo_mongodb}\end{figure}
